% Progress Report II — due Dec 1, 2 pages, 11 pt.
% Build: make report2
\documentclass[11pt]{article}
\usepackage{deqreport}
\reporttitle{Progress Report II}

\begin{document}
\makeheader

\section{Summary}

% Lead with the comparison result — that is the project's main question.
\textsc{TODO: one paragraph. The three backward methods now run on the same
convolutional model; the headline accuracy/time trade-off is X.}

\section{Progress Since Report I}

\subsection{Convolutional DEQ in PyTorch}

% Architecture, and what part of the backward pass is yours vs autograd.
The equilibrium layer is two $3\times3$ convolutions with group normalization
and input injection. The solver and the implicit gradient are implemented as a
custom \texttt{torch.autograd.Function}; autograd supplies only the
vector-Jacobian product of a \emph{single} layer evaluation, never a
differentiated solver trajectory.

\subsection{Three backward passes}

% The variants, stated as what replaces the inverse in Eq. (2) of Report I.
\begin{enumerate}
  \item \textbf{Exact implicit.} Solve $u = uJ + \dd{\ell}{\zstar}$ to
    tolerance with Anderson.
  \item \textbf{Truncated Neumann.} Approximate $(I-J)^{-1}$ by
    $\sum_{i=0}^{k} J^i$ for $k = 1 \ldots 5$~\cite{geng2021training}.
  \item \textbf{Jacobian-free.} Take $(I-J)^{-1} \approx I$, i.e.\
    $k = 0$~\cite{fung2022jfb}.
\end{enumerate}

\subsection{Unrolled baseline}

A weight-tied network applying the same layer $K$ times, trained with ordinary
backpropagation through all $K$ steps, gives the memory and accuracy reference
point.

\section{Results}

\begin{table}[h]
  \centering
  \caption{CIFAR-10, \textsc{TODO} epochs. Gradient cosine similarity is
    measured against the exact implicit gradient.}
  \begin{tabularx}{\textwidth}{@{}Xrrrr@{}}
    \toprule
    Method & Test acc.\ (\%) & Peak mem.\ (MB) & Time/epoch (s) & Grad.\ cos.\ sim. \\
    \midrule
    Unrolled ($K=\textsc{TODO}$) & \textsc{TODO} & \textsc{TODO} & \textsc{TODO} & -- \\
    Exact implicit & \textsc{TODO} & \textsc{TODO} & \textsc{TODO} & 1.000 \\
    Neumann ($k=1$) & \textsc{TODO} & \textsc{TODO} & \textsc{TODO} & \textsc{TODO} \\
    Neumann ($k=5$) & \textsc{TODO} & \textsc{TODO} & \textsc{TODO} & \textsc{TODO} \\
    Jacobian-free & \textsc{TODO} & \textsc{TODO} & \textsc{TODO} & \textsc{TODO} \\
    \bottomrule
  \end{tabularx}
  \label{tab:compare}
\end{table}

% A figure of solver iterations over training shows whether the model is
% becoming harder to solve — the main stability story.
% \begin{figure}[h]\centering
%   \includegraphics[width=0.75\textwidth]{figures/fwd_iters.pdf}
%   \caption{Forward solver iterations per epoch.}\label{fig:iters}
% \end{figure}

\textsc{TODO: two or three sentences reading the table.}

\section{Problems Encountered}

\textsc{TODO}

\section{Plan for Next Steps}

The final report (Dec 8) adds ablations on solver tolerance and Anderson
memory, the Neumann order sweep, Jacobian
regularization~\cite{bai2021stabilizing} if convergence degrades, and the full
write-up.

\printrefs
\end{document}
